\documentclass[11pt]{article}
\usepackage[margin=1in]{geometry}
\usepackage{amsmath,amssymb,amsthm,mathtools}
\usepackage{booktabs}
\usepackage[expansion=false]{microtype}
\usepackage{xcolor}
\usepackage[colorlinks=true,linkcolor=blue!50!black,citecolor=blue!50!black,urlcolor=blue!50!black]{hyperref}

\setlength{\emergencystretch}{3em}
\newtheorem{theorem}{Theorem}[section]
\newtheorem{proposition}[theorem]{Proposition}
\newtheorem{lemma}[theorem]{Lemma}
\newtheorem{corollary}[theorem]{Corollary}
\theoremstyle{definition}
\newtheorem{definition}[theorem]{Definition}
\newtheorem{example}[theorem]{Example}
\newtheorem{problem}[theorem]{Problem}
\theoremstyle{remark}
\newtheorem{remark}[theorem]{Remark}

\newcommand{\Z}{\mathbb Z}
\newcommand{\R}{\mathbb R}
\newcommand{\C}{\mathbb C}
\newcommand{\N}{\mathbb N}
\newcommand{\E}{\mathbb E}
\newcommand{\tr}{\operatorname{tr}}
\newcommand{\Ad}{\operatorname{Ad}}
\newcommand{\id}{\mathrm{id}}
\newcommand{\Var}{\operatorname{Var}}
\newcommand{\Cov}{\operatorname{Cov}}
\newcommand{\relu}{\operatorname{relu}}
\newcommand{\dep}[1]{|#1|}
\newcommand{\tail}{\mathrm{tail}}

\title{Pseudorandom arrows and PV cohomology\\[4pt]
\large small-bias diffusion on tail algebras, transverse heat on tiling cohomology,\\
and the K-theory of continuation (causal-state) hierarchies}
\author{Working note for the continuation-algebra programme\\
\small self-reviewed draft, not independently verified}
\date{3 October 2026}

\begin{document}
\maketitle

\begin{abstract}
We develop the expander and pseudorandomness side of the programme linking Cantor hierarchies, tail groupoids, convolution C$^*$-algebras and diffusion (Pearson--Bellissard, PB), with tiling theory (Bellissard--Benedetti--Gambaudo, BBG; Savinien--Bellissard, SB) and random ReLU networks as the intended application.

\emph{Arrows and expansion.} The binary tail groupoid is the orbit relation of $\Gamma=\bigoplus_{\N}\Z_2$, so its algebra carries a phase-space (Weyl) basis. In this basis (i) the PB heat semigroup lifted to the tail algebra is a mixed-unitary L\'evy semigroup whose multiplier depends only on the \emph{depth} of a Weyl label, and (ii) every average over a finite set of arrows is diagonal with eigenvalues equal to that set's symplectic Fourier biases. Consequences: $\varepsilon$-biased arrow sets are exactly $\varepsilon$-spectral sparsifiers of the conditional expectations. A sparse hierarchical L\'evy generator with $O(N^2/\varepsilon^2)$ jump unitaries (from one seed set of size $O(N/\varepsilon^2)$) generates a completely positive, trace-preserving quantum Markov semigroup whose Dirichlet form is $(1\pm\varepsilon)$-equivalent to the PB form and which commutes exactly with every refinement level. Heat-smoothed observables are fooled by states that are random only on the first $M(t,\varepsilon)=O(\log(1/t)+\log\log(1/\varepsilon))$ levels.

\emph{Neural (nontracial) states.} Admissible arrows of the continuation note are precisely the eigenoperators of the modular group. Hence its maximality theorem is an instance of Takesaki's theorem, and on each continuation class the state factorizes as a product. Dephasing moves stay state-preserving for free; transport moves must be Metropolized, with a distortion-dependent gap.

\emph{PV cohomology.} We read SB's PV complex as an arrow calculus. A BBG invariant measure turns it into a pre-Hilbert complex in which BBG's switching rules say exactly that the constant top cochain is harmonic, invariant measures are the positive functionals annihilating PV coboundaries, and gap labels are inner products with that harmonic constant. Transverse (PB-type) heat commutes with the PV coboundary at every level iff the hierarchy comes from an odometer (the fair binary case). Otherwise the commutator lives on the branch locus of the BBG approximants (the Kakutani--Rokhlin tower tops).

\emph{Continuation algebras.} Continuation classes are Crutchfield's causal states; the continuation algebra is the AF algebra of a ``causal-state Bratteli diagram''. KMS$_\beta$ states correspond to $\beta$-weighted switching rules, the BBG positive-cycle theorem in dimension zero. The neural state is the KMS$_1$ state given by the constant harmonic function. Finite rank gives at most $N$ extremal KMS states, and Birkhoff contraction gives uniqueness at a geometric rate. Simple finite-rank diagrams of rank $\ge2$ realize as one-dimensional FLC tilings with $H^1_{PV}\cong K_0$ of the continuation algebra. The golden-mean causal-state machine gives exactly the Fibonacci tiling.

\emph{Network bridges.} Partial annealing of discarded weight directions is Haar averaging over residual-rotation arrows, and it equals Gaussian heat flow on the activation. Downstream error is an exact positive spectral sum over moment orders. Depth acts on moment orders as a critical Galton--Watson process, with He initialization exactly at criticality.
\end{abstract}

\tableofcontents

\section{Status, conventions, sources}

\textbf{Status.} Items labelled Theorem, Proposition, Lemma or Corollary carry proofs here, with ``proof sketch'' marking where a step is delegated to a cited theorem. Finite-dimensional identities were also checked numerically (Appendix~\ref{app:num}; scripts in \texttt{code/}). The draft is self-reviewed only. Where a statement is known we say so; where we believe a formulation is new we say ``we have not found this stated'', which is weaker than claiming novelty.

\textbf{Sources.} PB: Pearson--Bellissard, \emph{Noncommutative Riemannian geometry and diffusion on ultrametric Cantor sets}, arXiv:0802.1336. BBG: Bellissard--Benedetti--Gambaudo, \emph{Spaces of tilings, finite telescopic approximations and gap-labelling}, arXiv:math/0109062 (Thm.~1.15, Def.~3.12, Prop.~5.3, 5.5, Thm.~5.6, Cor.~5.7, Thm.~1.16). SB: Savinien--Bellissard, \emph{A spectral sequence for the K-theory of tiling spaces}, arXiv:0705.2483 (Thm.~1, 2, 7; Def.~13--16; Prop.~3, 4; eq.~(1a),(1b)). PY: Proietti--Yamashita, arXiv:2006.08028. Other standard results are cited where used.

\textbf{The continuation note.} ``The continuation note'' is the user's nine-page note (continuation laws $\nu_p$, relation $p\sim_n q\iff\nu_p=\nu_q$, blocks $B_n=\bigoplus_{C\in P_n}M_{|C|}$, nested state-preserving expectations $F_n$, heat $T_t=(1-e^{-t\lambda_1})F_0+\sum_{n\ge1}(e^{-t\lambda_n}-e^{-t\lambda_{n+1}})F_n$, and the relaxation bound $D(\nu\|\tilde\nu)=I_\nu(P;Z\mid C)$). We use only these stated facts.

\section{Phase space of the binary tail algebra}\label{sec:phase}

Let $X_N=\{0,1\}^N$, $A_N=M_2^{\otimes N}\cong M_{2^N}$ with normalized trace $\tau$, and $A=\varinjlim A_N$ (via $a\mapsto a\otimes 1$), the UHF($2^\infty$) algebra $=C^*(R)$ of the tail relation $R$ on $X=\{0,1\}^{\N}$. Write $X,Z$ for the Pauli matrices.

\begin{definition}
The \emph{phase space} is $G_N=(\Z_2\times\Z_2)^N\ni w=(a,b)$. Its Weyl operators are $W(a,b)=\bigotimes_{i=1}^N X^{a_i}Z^{b_i}$, and its symplectic form is $\omega\big((a,b),(a',b')\big)=a\cdot b'+b\cdot a'\pmod 2$. The \emph{depth} $\dep w=\max\{i:(a_i,b_i)\neq(0,0)\}$ ($\dep 0=0$). For $0\le n\le N$, $G_{(n,N]}$ denotes the labels supported in coordinates $n+1,\dots,N$.
\end{definition}

\begin{proposition}[standard]\label{prop:weyl}
(i) $\{W(w)\}_{w\in G_N}$ is an orthonormal basis of $L^2(A_N,\tau)$ and $W(a,b)^*=(-1)^{a\cdot b}W(a,b)$. (ii) $W(g)W(w)W(g)^*=(-1)^{\omega(g,w)}W(w)$. (iii) Arrow reading: $W(a,0)=\sum_{q\in X_N}e_{q+a,q}$ is the full bisection ``translate by $a$'' of the tail groupoid; $W(0,b)$ is the character $\chi_b(\xi)=(-1)^{b\cdot\xi}$ in the Cartan subalgebra $C(X_N)$. The tail relation on $X$ is the orbit relation of the free translation action of $\Gamma=\bigoplus_{\N}\Z_2$, so $C^*(R)\cong C(X)\rtimes\Gamma$.
\end{proposition}

So a Weyl label is ``an arrow (translation) times a phase (character)'', and depth is the finest resolution it touches. This is the noncommutative version of the Walsh depth $|\omega|=\max\{i:\omega_i=1\}$ that governs the PB spectrum.

\section{The hierarchical heat as a L\'evy semigroup on phase space}\label{sec:levy}

Let $0=\lambda_0<\lambda_1<\dots$, $p_n(t)=e^{-t\lambda_n}-e^{-t\lambda_{n+1}}$ ($n<N$), $p_N(t)=e^{-t\lambda_N}$ (truncation at $N$; $p_0$ includes the $1-e^{-t\lambda_1}$ convention). Let $E_n:A_N\to A_n\otimes 1$ be the normalized partial trace and $T_t=\sum_{n=0}^N p_n(t)E_n$. For the canonical triadic PB example, $\lambda_m=2\sum_{j=0}^{m-2}9^j+4\cdot 9^{m-1}=\tfrac14(17\cdot 9^{m-1}-1)$, i.e.\ $4,38,344,3098,\dots$, so $\lambda_{m+1}=9\lambda_m+2$ for $m\ge1$.

\begin{theorem}[phase-space form]\label{thm:levy}
(a) $E_n=|G_{(n,N]}|^{-1}\sum_{g\in G_{(n,N]}}\Ad W(g)$.
(b) $T_tW(w)=e^{-t\lambda_{\dep w}}W(w)$ for every $w\in G_N$.
(c) $T_t=\sum_{g\in G_N}k_t(g)\Ad W(g)$, where $k_t=\sum_np_n(t)\,u_{(n,N]}$ and $u_H$ is the uniform probability on $H$. The symplectic Fourier transform is $\hat k_t(w):=\sum_gk_t(g)(-1)^{\omega(g,w)}=e^{-t\lambda_{\dep w}}$, and $k_t*k_s=k_{t+s}$. Thus $T_t$ is a mixed-unitary channel and $(k_t)$ is a convolution semigroup of probability measures on the compact group $G_N$ (and on $(\Z_2^2)^{\N}$ as $N\to\infty$). Its L\'evy measure is $J=\sum_{n<N}(\lambda_{n+1}-\lambda_n)u_{(n,N]}$: \emph{at rate $\lambda_{n+1}-\lambda_n$, re-randomize the whole phase space beyond level $n$.}
(d) On $C(X_N)$, $T_t$ is convolution by the $a$-marginal of $k_t$, i.e.\ the PB heat semigroup (PB Thm.~8/9 for the canonical spectrum).
\end{theorem}

\begin{proof}
(a) By Prop.~\ref{prop:weyl}(ii) the average multiplies $W(w)$ by $\E_g(-1)^{\omega(g,w)}$. That average is $1$ if $w|_{(n,N]}=0$ and $0$ otherwise, since $\omega$ restricted to $G_{(n,N]}$ is nondegenerate. The normalized partial trace acts the same way, because the trace of a nontrivial Pauli string vanishes. (b) $T_tW(w)=\sum_{n\ge\dep w}p_n W(w)=e^{-t\lambda_{\dep w}}W(w)$. (c) Combine (a) and (b). The characters $g\mapsto(-1)^{\omega(g,w)}$ exhaust $\widehat{G_N}$, so the convolution identity follows from $\hat k_t\hat k_s=\hat k_{t+s}$. (d) $\Ad W(a,b)$ acts on $C(X_N)$ as translation by $a$.
\end{proof}

\begin{remark}
The diagonal statement belongs to the theory of isotropic Markov semigroups on ultrametric spaces (Bendikov--Grigor'yan--Pittet--Woess). The point here is the noncommutative, mixed-unitary, Weyl-diagonal form. In it the PB diffusion becomes a random walk on the arrow phase space, and that walk is what we sparsify.
\end{remark}

\section{Arrow averages are exactly their biases}\label{sec:bias}

\begin{definition}
For a probability $\pi$ on $G_N$, the \emph{arrow average} is $\Phi_\pi=\sum_g\pi(g)\Ad W(g)$ and its \emph{bias} is $\hat\pi(w)=\sum_g\pi(g)(-1)^{\omega(g,w)}$. If $\pi$ is supported in $G_{(n,N]}$, it is \emph{$\varepsilon$-biased at level $n$} if $|\hat\pi(w)|\le\varepsilon$ whenever $w|_{(n,N]}\neq0$.
\end{definition}

\begin{lemma}\label{lem:bias}
$\Phi_\pi W(w)=\hat\pi(w)W(w)$. Moreover $\Phi_\pi$ is unital, completely positive, $\tau$-preserving, and $\tau$-symmetric when $\pi$ is symmetric (always the case here, since $\Ad W(g)$ is self-adjoint on $L^2(\tau)$).
\end{lemma}

\begin{proposition}[bias $=$ sparsification]\label{prop:sparsify}
Let $\pi$ be supported in $G_{(n,N]}$ with $\varepsilon:=\max_{w|_{(n,N]}\ne0}|\hat\pi(w)|$. Then $\Phi_\pi|_{A_n}=\id$ and $\|\Phi_\pi^k-E_n\|_{L^2\to L^2}=\varepsilon^k$. Noncommutative expander mixing: $|\tau(y^*\Phi_\pi x)-\tau(y^*E_nx)|\le\varepsilon\|x-E_nx\|_2\|y-E_ny\|_2$. In particular, for projections $P,Q$ and $n=0$,
\[|\tau(P\Phi_\pi Q)-\tau(P)\tau(Q)|\le\varepsilon\sqrt{\tau(P)(1-\tau(P))\tau(Q)(1-\tau(Q))}.\]
On $C(X_N)$ the eigenvalues are the ordinary biases of the $a$-marginal, which is the Alon--Roichman equivalence ``Cayley graph on $\Z_2^m$ is an $\varepsilon$-expander $\iff$ generators are $\varepsilon$-biased''.
\end{proposition}

\begin{proof}
Everything is Weyl-diagonal. $E_n$ is the projection onto labels with $w|_{(n,N]}=0$, on which $\hat\pi=1$; on the complement $|\hat\pi|\le\varepsilon$. For the projection estimate use $\|P-\tau(P)\|_2^2=\tau(P)(1-\tau(P))$.
\end{proof}

\begin{remark}
For $n=0$ this is the mechanism behind Ambainis--Smith's small-bias Pauli sets (approximate quantum encryption) and the abelian case of quantum expanders. What follows uses it hierarchically, level by level, inside a diffusion.
\end{remark}

\section{Sparse hierarchical sparsifiers of the PB diffusion}\label{sec:sparse}

\begin{theorem}[sparse hierarchical L\'evy sparsifier]\label{thm:sparse}
For $0\le n<N$ let $\pi_n$ be supported in $G_{(n,N]}$ and $\varepsilon_n$-biased at level $n$, with $\varepsilon=\max_n\varepsilon_n<1$. Put
\[\hat L=\sum_{n<N}(\lambda_{n+1}-\lambda_n)\big(\Phi_{\pi_n}-\id\big),\qquad \hat T_t=e^{t\hat L}.\]
Then:
\begin{enumerate}
\item $(\hat T_t)$ is a semigroup of unital, completely positive, trace-preserving, $\tau$-symmetric, mixed-unitary maps, with $\hat T_tW(w)=e^{-t\hat\psi(w)}W(w)$ and $\hat\psi(w)=\sum_{n<\dep w}(\lambda_{n+1}-\lambda_n)(1-\hat\pi_n(w))$.
\item \emph{Two-sided spectral equivalence:} $(1-\varepsilon)\lambda_{\dep w}\le\hat\psi(w)\le(1+\varepsilon)\lambda_{\dep w}$ for all $w$. Hence $(1-\varepsilon)Q\le\hat Q\le(1+\varepsilon)Q$ for the Dirichlet forms $Q(x)=-\tau(x^*Lx)=\sum_w\lambda_{\dep w}|\hat x(w)|^2$ and $\hat Q$.
\item \emph{Gap:} $\|\hat T_tx-\tau(x)\|_2\le e^{-(1-\varepsilon)\lambda_1t}\|x-\tau(x)\|_2$.
\item \emph{Uniform closeness:} $\sup_{t\ge0}\|\hat T_t-T_t\|_{L^2\to L^2}\le\varepsilon/(e(1-\varepsilon))$.
\item \emph{Exact refinement compatibility:} $\hat T_tE_m=E_m\hat T_t$ and $\hat T_t(A_m)\subseteq A_m$ for every $m$.
\item \emph{Classical restriction:} on $C(X_N)$, $\hat T_t$ is the Markov jump process $\xi\mapsto\xi+a$ at rate $\sum_n(\lambda_{n+1}-\lambda_n)\pi_n^X(a)$, a spectral $(1\pm\varepsilon)$-sparsifier of the PB Laplacian.
\end{enumerate}
\end{theorem}

\begin{proof}
(1) $\hat L(x)=\sum_j\gamma_j(V_j^*xV_j-x)$ with unitary $V_j$ and $\gamma_j\ge0$, a Lindblad generator with zero Hamiltonian, so it generates a CP unital semigroup. Each $\Ad W(g)$ preserves $\tau$ and is $\tau$-symmetric. The Poisson expansion of $e^{t\hat L}$ is a convex combination of products of $\Ad W(g)$'s. The multiplier follows from Lemma~\ref{lem:bias}, using $\hat\pi_n(w)=1$ for $n\ge\dep w$, because then $w|_{(n,N]}=0$.
(2) For $n<\dep w$, $w|_{(n,N]}\ne0$, so $1-\hat\pi_n(w)\in[1-\varepsilon,1+\varepsilon]$; then sum, using $\sum_{n<m}(\lambda_{n+1}-\lambda_n)=\lambda_m$.
(3) Follows from (2) with $\dep w\ge1$.
(4) With $x=t\lambda$ and $\theta\in[1-\varepsilon,1+\varepsilon]$: $e^{-(1-\varepsilon)x}-e^{-x}\le\varepsilon xe^{-(1-\varepsilon)x}\le\varepsilon/(e(1-\varepsilon))$ and $e^{-x}-e^{-(1+\varepsilon)x}\le\varepsilon xe^{-x}\le\varepsilon/e$.
(5) $\hat T_t$ and $E_m$ are both Weyl-diagonal, and $\hat T_t$ does not increase depth.
(6) For labels $(0,b)$, $\omega((a,b'),(0,b))=a\cdot b$.
\end{proof}

\begin{proposition}[one seed for all levels; sizes]\label{prop:seed}
If $S\subset G_N$ is $\varepsilon$-biased on all of $G_N$, the truncations $\pi_n=\text{law of }g|_{(n,N]}$ ($g\sim\mathrm{Unif}(S)$) are $\varepsilon$-biased at level $n$, because $\omega(g|_{(n,N]},w)=\omega(g,w)$ for $w$ supported in $(n,N]$. A uniformly random $S$ with $|S|\ge 2\varepsilon^{-2}(2N\ln2+\ln(2/\delta))$ is $\varepsilon$-biased with probability $\ge1-\delta$ (Hoeffding plus a union bound over $4^N$ labels). Explicit constructions (Alon--Goldreich--H\aa stad--Peralta) give $|S|=O((N/\varepsilon)^2)$. Hence the sparse diffusion uses at most $N|S|=O(N^2/\varepsilon^2)$ distinct jump unitaries, generated by one seed set of size $O(N/\varepsilon^2)$. Each jump consumes $\log_2|S|=\log_2N+2\log_2(1/\varepsilon)+O(1)$ random bits instead of the $2(N-n)$ bits an exact level-$n$ re-randomization needs.
\end{proposition}

\begin{proposition}[necessity]\label{prop:lower}
If $L=\sum_gc_g(\Ad W(g)-\id)$ with $c\ge0$ has $\psi(w)>0$ for all $w\ne0$, then $\mathrm{supp}(c)$ spans $G_N$. Hence at least $2N$ jump unitaries are necessary.
\end{proposition}

\begin{proof}
$\psi(w)=2\sum_{g:\,\omega(g,w)=1}c_g$. If the support spans a proper subspace $V$, nondegeneracy of $\omega$ gives $w\ne0$ with $\omega(V,w)=0$, so $\psi(w)=0$.
\end{proof}

\begin{remark}
This is the hierarchical analogue of ``expanders sparsify complete graphs'' (Spielman--Teng, Batson--Spielman--Srivastava), here for the ultrametric PB Laplacian and its quantum lift, with the semigroup law and exact refinement compatibility kept. The gap between $O(N^2/\varepsilon^2)$ and $\Omega(N)$ distinct jumps is open.
\end{remark}

\section{Hierarchy-adapted pseudorandomness}\label{sec:seed}

\begin{theorem}[seed length for heat-smoothed observables]\label{thm:seedlen}
Let $\rho\ge0$ with $\tau(\rho)=1$ and put $\hat\rho(w)=\tau(\rho W(w))$, so $|\hat\rho(w)|\le1$. For every $x\in A_N$,
\[|\tau(\rho\,T_tx)-\tau(x)|\le\|x-\tau(x)\|_2\Big(\sum_{w\ne0}e^{-2t\lambda_{\dep w}}|\hat\rho(w)|^2\Big)^{1/2}.\]
Let $c_m=\#\{w:\dep w=m\}=3\cdot4^{m-1}$ (classically, $2^{m-1}$ characters). If $|\hat\rho(w)|\le\delta$ for $0<\dep w\le M$ (e.g.\ $\delta=0$ when the marginal of $\rho$ on the first $M$ qubits is maximally mixed), then the error is at most
\[\|x-\tau(x)\|_2\Big(\delta^2\textstyle\sum_{m\le M}c_me^{-2t\lambda_m}+\sum_{m>M}c_me^{-2t\lambda_m}\Big)^{1/2}.\]
For the canonical spectrum, $\lambda_m\ge4\cdot9^{m-1}$, and the tail is $\le\varepsilon^2$ as soon as $64t\,9^M\ge\ln8$ and $8t\,9^M\ge2\ln(1/\varepsilon)+M\ln4+\ln6$. Hence the needed true randomness is $M(t,\varepsilon)=O(\log(1/t)+\log\log(1/\varepsilon))$ levels, \emph{independent of $N$}.
\end{theorem}

\begin{proof}
Expand $x$ and $\rho$ in the Weyl basis, use Theorem~\ref{thm:levy}(b) and Cauchy--Schwarz, and note that the $w=0$ term is $\tau(x)$. For the tail, the first condition makes consecutive terms decrease by a factor $\le1/2$, so the tail is at most twice its first term, $2\cdot3\cdot4^Me^{-8t9^M}$.
\end{proof}

\begin{center}
\begin{tabular}{lcccc}
\toprule
$M(t,\varepsilon)$ (classical $=$ noncommutative here) & $\varepsilon=10^{-3}$ & $10^{-6}$ & $10^{-12}$ & $10^{-24}$\\
\midrule
$t=10^{-4}$ & 5 & 5 & 6 & 6\\
$t=10^{-3}$ & 4 & 4 & 5 & 5\\
$t=10^{-2}$ & 3 & 3 & 3 & 4\\
$t=10^{-1}$ & 2 & 2 & 2 & 3\\
\bottomrule
\end{tabular}
\end{center}

\begin{corollary}[address completion]
Let an \emph{address} be a prefix $p$ of length $n$, with state $\rho_p=2^ne_p\otimes1$. Replace the tail by any $\tilde\rho=2^ne_p\otimes\sigma$ in which $\sigma$ is maximally mixed on qubits $n+1,\dots,M$ and arbitrary beyond. Then $|\tau(\rho_pT_tx)-\tau(\tilde\rho T_tx)|\le2\|x\|_2(\sum_{m>M}c_me^{-2t\lambda_m})^{1/2}$. If $n\ge M(t,\varepsilon)$, no randomness at all is needed: after diffusion time $t$, resolution beyond $M$ is invisible.
\end{corollary}

\begin{remark}
Contrast this with Bonami--Beckner noise, whose multiplier $\rho^{|\omega|_1}$ depends on Hamming weight: there the number of characters above a threshold grows polynomially in $N$ and fooling needs $k$-wise independence with $k\sim\log(1/\varepsilon)/\log(1/\rho)$. The PB hierarchy is \emph{pseudorandomly cheap} because its multipliers decay doubly exponentially in depth while label counts grow only exponentially.
\end{remark}

\section{Neural (nontracial) states: modular admissibility}\label{sec:modular}

Let $\nu$ be a probability on $X$ with continuation laws $\nu_p$ and density $w=d\nu/d\mu$ relative to the fair measure, assumed positive. (Zero-gain support is treated by restriction, as in the continuation note.) Its state on the tail algebra has modular group $\sigma_t(f)(\xi,\zeta)=(w(\xi)/w(\zeta))^{it}f(\xi,\zeta)$ (Renault); in finite truncation, $\sigma_t=\Ad\rho^{it}$ with $\rho=w\in C(X_N)$.

\begin{theorem}[admissible arrows are modular eigenoperators]\label{thm:modular}
For prefixes $p,q$ of length $n$, the arrow bisection $e_{pq}\otimes1$ satisfies $\sigma_t(e_{pq}\otimes1)=\lambda^{it}e_{pq}\otimes1$ for all $t$ and some $\lambda>0$ iff $\nu_p=\nu_q$. In that case $\lambda=\nu[p]/\nu[q]$.
\end{theorem}

\begin{proof}
$w(pz)=2^n\nu[p]g_p(z)$ with $g_p=d\nu_p/d\mu_\tail$. So $\sigma_t(e_{pq}\otimes1)=(\nu[p]/\nu[q])^{it}\,e_{pq}\otimes M_{(g_p/g_q)^{it}}$. This is a scalar multiple of $e_{pq}\otimes1$ iff $g_p/g_q$ is a.e.\ constant, and since both are probability densities, iff $g_p=g_q$.
\end{proof}

\begin{corollary}[the maximality theorem via Takesaki]\label{cor:takesaki}
Let $D_n\subseteq\mathcal B\subseteq A_n\otimes1$ with $D_n$ the level-$n$ diagonal. Then $\mathcal B$ is $\sigma$-invariant iff $\mathcal B\subseteq B_n=\mathrm{span}\{e_{pq}:p\sim_nq\}$, and $B_n$ is $\sigma$-invariant. By Takesaki's theorem (a $\varphi$-preserving normal conditional expectation onto a von Neumann subalgebra exists iff the subalgebra is $\sigma^\varphi$-invariant), $B_n$ is the largest such subalgebra admitting a state-preserving expectation. This recovers the maximality statement of the continuation note.
\end{corollary}

\begin{proof}
A $*$-subalgebra containing $D_n$ is spanned by the matrix units it contains, obtained by compressing with diagonal projections. If $e_{pq}\in\mathcal B$ is not admissible, then $\sigma_t(e_{pq})\notin A_n\otimes1$ by the formula above.
\end{proof}

\begin{proposition}[product factorization]\label{prop:product}
On the corner of a class $C\in P_n$, $e_CAe_C\cong M_{|C|}\otimes A_{>n}$, the normalized state is $\rho_C\otimes\varphi_{\nu_C}$ with $\rho_C=\mathrm{diag}(\nu[p]/\nu[C])_{p\in C}$. Conversely, if the corner state is a product with diagonal first factor, all $p\in C$ have the same continuation law. Admissibility is thus exactly ``the neural weight splits as prefix mass $\otimes$ shared continuation''.
\end{proposition}

\begin{proposition}[which pseudorandom moves survive a neural state]\label{prop:moves}
(a) \emph{Dephasing is free.} Every $\Ad W(0,b)$ commutes with the diagonal density, so it is state-preserving and KMS-symmetric. Small-bias phase sets sparsify the expectation onto $A_n\otimes C(X_{(n,N]})$ exactly as in Section~\ref{sec:sparse}, and that expectation is state-preserving because the density lies in its range.
(b) \emph{Transport must be Metropolized.} For a prefix $p$ with continuation density $g=g_p>0$ on $X_{(n,N]}$ and $S$ $\varepsilon$-biased, the chain $Kf(z)=f(z)+|S|^{-1}\sum_{s\in S}\min(1,g(z+s)/g(z))(f(z+s)-f(z))$ is $g\mu$-reversible and has spectral gap $\ge(1-\varepsilon)\min g/\max g$. If $p\sim q$, the chains on $[p]$ and $[q]$ coincide, so $\bigoplus_p\id_p\otimes K_{g_p}$ commutes with all admissible arrows of $B_n$.
\end{proposition}

\begin{proof}
(b) Detailed balance holds since $\mu(z)\min(g(z),g(z+s))$ is symmetric. The Dirichlet form is $\ge(\min g)\,\mathcal E_S$, and $\mathcal E_S\ge(1-\varepsilon)\Var_\mu$ by Proposition~\ref{prop:sparsify}. Finally $\Var_{g\mu}(f)\le\int(f-\mu f)^2g\,d\mu\le\max g\,\Var_\mu(f)$.
\end{proof}

\section{The shift semigroupoid renormalizes the heat}\label{sec:shift}

The one-sided shift semigroupoid $\{(x,k,l,y):\sigma^kx=\sigma^ly\}$ has Cuntz algebra $\mathcal O_2$, whose gauge-fixed core is the tail algebra. The canonical endomorphism $\iota(x)=s_0xs_0^*+s_1xs_1^*=1\otimes x$ (``prepend a free symbol'') is the semigroupoid arrow acting on the core.

\begin{proposition}[renormalization identity]\label{prop:rg}
For the canonical triadic spectrum,
\[T_t\circ\iota=\iota\circ\big(e^{-2t}T_{9t}+(1-e^{-2t})E_0\big).\]
\end{proposition}

\begin{proof}
$\iota W(w)=W(0\oplus w)$ has depth $\dep w+1$, and $\lambda_{m+1}=9\lambda_m+2$ for $m\ge1$, while $\iota(1)=1$.
\end{proof}

Moving one step up the semigroupoid dilates time by $9=3^2$ (the diffusive scaling of the triadic metric) and adds depolarization at rate $2$. This is the exact form, for the PB hierarchy, of the level shift that substitution (Anderson--Putnam/BBG inflation) induces on tiling hierarchies.

\section{PV cohomology through arrows}\label{sec:pv}

\subsection{What SB and BBG provide}
\emph{SB.} For an aperiodic repetitive FLC tiling of $\R^d$ with punctured $\Delta$-complex structure, the $\Delta$-transversal $\Xi_\Delta=\bigsqcup_n\Xi^n_\Delta$ is a Cantor set (Def.~13). The operators $\theta_{\sigma\tau}=\chi_\sigma t_{x_{\sigma\tau}}\chi_\tau$ are partial isometries implementing the generating arrows of the $\Delta$-transversal groupoid (Def.~14, Rem.~3), with relations (1a) $\theta\theta^*=\chi_\sigma$ and (1b) $\sum_{\sigma\supset\tau}\theta^*_{\sigma\tau}\theta_{\sigma\tau}=\chi_\tau$. The PV complex is $C^n_{PV}=C(\Xi^n_\Delta,\Z)$ with $d^n_{PV}=\sum_{\sigma\in S^n_0}\sum_i(-1)^i\theta_{\sigma\partial_i\sigma}$ (Def.~15). On functions constant on level-$l$ acceptance zones, $d_{PV}$ is the simplicial coboundary of the patch space $B_l$ (Prop.~3, using Lemma~1: the zone of $\sigma$ is partitioned by the zones of its preimages). PV is the direct limit over levels (Prop.~4), hence isomorphic to \v{C}ech cohomology of the hull (Thm.~7). A spectral sequence $E_2^{rs}=H^r_{PV}(B_0;K^s(\Xi_\Delta))\Rightarrow K_{r+s+d}(C(\Omega)\rtimes\R^d)$ follows (Thm.~1). For $\Z^d$-actions the PV differential is $\sum_i(\alpha_{i*}-1)\otimes e_i\wedge$ (Thm.~2).

\emph{BBG.} The hull is an inverse limit of an expanding flattening sequence of branched oriented flat manifolds (Thm.~1.15). Invariant measures $\cong$ coherent families of positive weights on top cells satisfying the switching (Kirchhoff) rules, i.e.\ positive $d$-cycles (Def.~3.12, Prop.~5.5, Thm.~5.6). Bounded $\dim H_d$ gives at most $N$ ergodic measures, and bounded incidence coefficients give unique ergodicity (Cor.~5.7). The gap-labelling theorem states $\mathcal T_\mu(K_0)=\int d\mu_t\,C(\Gamma,\Z)$ (Thm.~1.16).

\emph{PY.} Groupoid homology spectral sequence $E^2_{pq}=H_p(G,K_q)\Rightarrow K_{p+q}(C^*_rG)$ for torsion-free ample groupoids with strong Baum--Connes. For tiling transversals, $H^k(G|_X)\cong H_{d-k}(G|_X)$ by Poincar\'e duality for $\Z^d$ (PY \S4.2). Comparing this with the SB sequence is stated as open (PY Rem.~4.1).

\subsection{BBG measures make PV a Hodge complex}

Let $F_l^n=C(\Sigma^n_l)$ be the level-$l$ PV cochains (SB Def.~16), and let $\mu$ be a finite measure on $\Xi_\Delta$ invariant under the $\Delta$-transversal groupoid.

\begin{theorem}[PV Hodge structure from BBG measures]\label{thm:pvhodge}
\begin{enumerate}
\item Each $\theta_{\sigma\tau}$ is a partial isometry on $L^2(\Xi_\Delta,\mu)$, with $\theta^*\theta$ the indicator of $\{\xi\in\Xi(\tau):t_x\xi\in\Xi(\sigma)\}$. So the PV ring is represented by partial isometries and $d_{PV}$ is a bounded operator with an adjoint $d^*_\mu$.
\item On $F_l$ the $L^2(\mu)$ inner product is the weighted simplicial inner product $\langle f,g\rangle=\sum_{\sigma_l}f(\sigma_l)g(\sigma_l)w_l(\sigma_l)$, with $w_l(\sigma_l)=\mu(\Xi_{l,\Delta}(\sigma_l))$, and the inclusions $F_l\subset F_{l+1}$ are isometric.
\item \emph{Switching rules $=$ harmonicity of $1$.} For positive weights $w$ on the $d$-simplices of $B_l$ (and any positive weights below), $d^*_w\mathbf 1=0$ iff $\sum_\sigma w(\sigma)\sigma$ is a simplicial $d$-cycle, i.e.\ iff BBG's switching rules hold. In particular, for invariant $\mu$ the constant top cochain is harmonic at every level.
\item \emph{PV form of BBG Thm.~5.6 (proof sketch).} A finite positive measure $\mu$ on $\Xi^d_\Delta$ is the restriction of an invariant measure iff $\mu(d_{PV}\phi)=0$ for all $\phi\in C(\Xi^{d-1}_\Delta,\Z)$. So invariant measures are the positive functionals on top PV cochains that annihilate PV coboundaries: a positive cone in $(H^d_{PV}\otimes\R)^*$.
\item \emph{Gap labels.} For $f\in C(\Xi^d_\Delta,\Z)$, $\mu(f)=\langle f,\mathbf1\rangle_\mu$. The $\Z$-module of these values is BBG's gap-label module (Thm.~1.16; zones of top simplices are translates of tile zones). The label of a class equals the inner product of its harmonic projection with the harmonic constant.
\end{enumerate}
\end{theorem}

\begin{proof}
(1) Invariance makes $t_x$ measure-preserving between the relevant zones, and the computation of $\theta^*\theta$ is direct. (2) Follows from SB Lemma~1 and additivity. (3) $\langle d\phi,\mathbf1\rangle_w=\sum_\sigma w(\sigma)\sum_i(-1)^i\phi(\partial_i\sigma)=\phi(\partial w)$ for every $(d-1)$-cochain $\phi$, independently of the lower weights.
(4) $\Rightarrow$: by (3) at every level. $\Leftarrow$: $C(\Xi^{d-1}_\Delta)=\bigcup_lF_l^{d-1}$ and $d$ preserves levels, so $\partial w_l=0$ for all $l$. Interior faces of a tile force equal weights on its simplices. The resulting coherent positive cycles give an invariant measure by BBG Thm.~5.6, whose argument transfers verbatim to proper sequences of patch spaces. It agrees with $\mu$ on all zones, which generate.
(5) Is (4) together with BBG Prop.~5.9 and Thm.~1.16.
\end{proof}

\begin{remark}[expansion lives transversally]
The leafwise structure carried by $d_{PV}$ has no spectral gap. The $\Delta$-transversal relation is hyperfinite (amenable), and ergodic hyperfinite p.m.p.\ relations are never strongly ergodic (Connes--Weiss, Schmidt, Jones--Schmidt), whereas a spectral gap for a finitely generated Markov operator would force strong ergodicity. So $d^*_\mu d_{PV}$ restricted to $L^2_0(\Xi^0_\Delta)$ is gapless. The transverse PB-type heat $\sum p_lE_l$ has gap $\lambda_1$ by construction. In this sense the PV complex couples a \emph{gapless leafwise} differential with a \emph{gapped transverse} hierarchy, and the expanders of Sections~\ref{sec:bias}--\ref{sec:seed} act on the transverse side.
\end{remark}

\subsection{When does transverse heat commute with the PV coboundary?}

Take $d=1$ in the model $(\Xi,\theta)$ (minimal Cantor system) with PV complex $0\to C(\Xi,\Z)\xrightarrow{1-U}C(\Xi,\Z)\to0$, $Uf=f\circ\theta$ (SB Thm.~2, \S4.3). Let $\mu$ be invariant with full support, let $(P_l)$ be nested clopen partitions generating the topology, and let $E_l$ be the $\mu$-conditional expectations.

\begin{theorem}[odometer characterization]\label{thm:odometer}
(a) $E_lU=UE_l$ iff $\theta$ permutes the atoms of $P_l$. By minimality the permutation is a single cycle, giving a factor map onto $\Z/|P_l|$.
(b) If this holds for all $l$, then $(\Xi,\theta)$ is conjugate to the odometer $\varprojlim\Z/|P_l|$ and the $P_l$ are the pulled-back cylinder partitions. In that case $T_t=\sum p_lE_l$ is a PV chain map that preserves gap labels ($\mu(T_tf)=\mu(f)$) and retracts top cochains exponentially onto their harmonic part: $\|T_tf-\mu(f)\|_2\le e^{-t\lambda_1}\|f-\mu(f)\|_2$.
(c) If $(\Xi,\theta)$ has no nontrivial finite factor (Sturmian, weakly mixing), no $E_l$ with $|P_l|>1$ commutes with $U$.
\end{theorem}

\begin{proof}
(a) A unitary commutes with a finite-rank orthogonal projection iff it maps the range onto itself. For $U$ this means $f\circ\theta$ is $P_l$-measurable for every $P_l$-measurable $f$, i.e.\ $\theta^{-1}$ maps atoms to unions of atoms, and counting forces single atoms. A union of an orbit of atoms is clopen and invariant, so minimality gives a single cycle. (b) The factor maps separate points because the partitions generate. Since each $E_l$ commutes with $U$, $T_t$ commutes with $1-U$; it is $\mu$-preserving, and its gap is $\lambda_1$ as in Theorem~\ref{thm:levy}. (c) An invariant nontrivial finite partition is a finite factor.
\end{proof}

So the fair binary PB hierarchy is precisely the dyadic-odometer case, where diffusion and PV are perfectly compatible. Every genuinely aperiodic FLC tiling sits off this case, and the defect is quantified next.

\begin{proposition}[the defect lives on the branch locus]\label{prop:kr}
Let $P_l$ be Kakutani--Rokhlin partitions (towers $\{B_i,\theta B_i,\dots,\theta^{h_i-1}B_i\}$, the $d=1$ BBG approximant being the branched graph whose loops are the towers). Then $E_lU=UE_l$ off the tower tops, and
\[\|(E_lU-UE_l)f\|_2\le2\|f\|_\infty\,\mu\Big(\bigcup_{i\in I^*}\theta^{h_i-1}B_i\Big)^{1/2},\]
where $I^*$ are the towers whose top is mapped into at least two base atoms. Hence $\|[T_t,U]f\|_2\le2\|f\|_\infty\sum_lp_l(t)\,\mu(\text{branching tops at level }l)^{1/2}$. A single tower (an odometer level) has no defect.
\end{proposition}

\begin{proof}
$\theta$ maps the non-top level $\theta^jB_i$ measure-preservingly onto the atom $\theta^{j+1}B_i$, so both sides agree there. A top mapped into a single base atom also agrees. On the remaining tops the two averages differ by at most $2\|f\|_\infty$.
\end{proof}

\begin{problem}
Extend Theorem~\ref{thm:odometer} and Proposition~\ref{prop:kr} to $d\ge2$. Conjecture: commutation at all levels $\iff$ limit-periodic hierarchy. The defect of $[E_l,d_{PV}]$ should be bounded by the BBG transverse weight of the branch locus of $B_l$, via the $L^2$ version of SB eq.~(5).
\end{problem}

\section{Continuation algebras as causal-state Bratteli diagrams}\label{sec:causal}

\subsection{The diagram}
Restrict to positive-mass prefixes. The \emph{continuation diagram} has vertices $V_n=P_n$ (continuation classes) and one edge $(C,b)$ from $C$ to $C\cdot b$ for each symbol $b$ with $\nu_C(b)>0$; this is well defined because $p\sim q\Rightarrow pb\sim qb$. Paths from the root to $C$ are the prefixes in $C$, and infinite paths are $\mathrm{supp}\,\nu$. Tail equivalence of paths is the continuation groupoid $G_B$: $(pz,qz)$ with $p\sim q$. Hence:

\begin{proposition}
$B=\varinjlim B_n$ is the AF algebra of the continuation diagram, with $K_0(B)=\varinjlim(\Z^{P_n},A_n)$, $A_n(C',C)=\#\{b:C\cdot b=C'\}$. By Matui's theorem for AF groupoids, $H_0(G_B;\Z)\cong K_0(B)$ and $H_{k}(G_B)=0$ for $k>0$.
\end{proposition}

\begin{remark}[causal states]
$p\sim q\iff\nu_p=\nu_q$ is exactly Crutchfield--Young's causal-state equivalence of histories (finite-history, time-inhomogeneous form), and the refinement property is unifilarity. So the continuation algebra is the C$^*$-algebra of an $\varepsilon$-machine. The relaxation cost $I_\nu(P;Z\mid C)$ of the continuation note is the predictive information discarded by coarse-graining histories, the quantity optimized in causal-state information bottlenecks (Still--Crutchfield--Ellison; Creutzig--Globerson--Tishby). For generic neural gains no two prefixes have proportional continuation profiles, so all classes are singletons, $B=C(X)$ and $K_0=C(X,\Z)$: infinite statistical complexity. Structure appears only after relaxation.
\end{remark}

\subsection{KMS states are weighted positive cycles}

Let $\alpha_t(e_{pq})=(\nu[q]/\nu[p])^{it}e_{pq}$. This is consistent across levels because $\nu[pb]/\nu[qb]=\nu[p]/\nu[q]$ for $p\sim q$; it is the modular dynamics of $\nu$, with time reversed to make $\nu$ KMS$_1$.

\begin{theorem}[KMS$_\beta$ $\leftrightarrow$ $\beta$-weighted switching rules]\label{thm:kms}
KMS$_\beta$ states of $(B,\alpha)$ correspond bijectively to families $x^{(n)}_C\ge0$ with $x^{(0)}_\emptyset=1$ and
\[x^{(n)}_C=\sum_{b:\,\nu_C(b)>0}\nu_C(b)^\beta\,x^{(n+1)}_{C\cdot b},\]
via $\psi(e_{pq})=\delta_{pq}\,x_C\,\nu[p]^\beta$ for $p\in C$. The set is a nonempty compact simplex for every $\beta$.
\begin{itemize}
\item $\beta=0$: traces, i.e.\ unweighted switching rules. This is the zero-dimensional form of BBG Thm.~5.6, and gap labels are the $\Z$-module of class masses.
\item $\beta=1$: $x\equiv1$ is the neural state $\nu$. In general KMS$_1$ states are the positive space--time harmonic functions $h(n,C)$ of the \emph{class chain} $C_n\to C_n\cdot b$ (probability $\nu_{C_n}(b)$), normalized by $h(0,\emptyset)=1$. $\nu$ is extremal iff the invariant $\sigma$-field $\bigcap_n\sigma(C_n,\xi_{>n})$ is $\nu$-trivial, and $\nu$ is the unique KMS$_1$ state iff the minimal Martin boundary of the class chain is a point.
\end{itemize}
\end{theorem}

\begin{proof}
$\alpha$ leaves each $B_n=\bigoplus_CM_C$ invariant, with $\alpha_t=\Ad e^{-itH_n}$, $H_n=\mathrm{diag}(\log\nu[p])$. KMS$_\beta$ states on a direct sum of matrix algebras are convex combinations of block Gibbs states $\propto\tr(e^{\beta H}\cdot)$, which gives $\psi(e_{pp})=x_C\nu[p]^\beta$. Consistency $\psi(e_{pp})=\sum_b\psi(e_{pb,pb})$ together with $\nu[pb]=\nu[p]\nu_C(b)$ gives the recursion. KMS on the dense analytic subalgebra $\bigcup B_n$ suffices. Nonemptiness: the level-$\le m$ solution sets are nonempty, compact (each $x_C$ is bounded by the reciprocal of a path weight) and nested, so the projective limit is nonempty. For $\beta=1$, $\psi=\nu^h$ with density $h(n,C_n)$, a positive martingale. Extremality and ergodicity follow by the standard Doob/ergodic-decomposition argument.
\end{proof}

\begin{corollary}[finite rank, as in BBG Cor.~5.7]
If $|P_n|\le N$ for all $n$, then for each $\beta$ there are at most $N$ extremal KMS$_\beta$ states. The proof is BBG's: a nested intersection of polytopes with $\le N$ vertices.
\end{corollary}

\begin{proposition}[Birkhoff contraction $\Rightarrow$ uniqueness, with rate]\label{prop:birkhoff}
Let $M_n(\beta)$ be the $P_n\times P_{n+1}$ weighted incidence matrices. If, after telescoping, infinitely many blocks are strictly positive with projective diameter $\le D$, the KMS$_\beta$ state is unique. The level-$n$ cones contract in Hilbert's projective metric by $\tanh(D/4)$ per block (Birkhoff--Hopf). For $\beta=1$ this is Dobrushin/Birkhoff mixing of the class chain: \emph{fast mixing of the continuation hierarchy gives uniqueness of the neural KMS state, and deep continuation information is forgotten geometrically.}
\end{proposition}

\begin{example}[the golden-mean $\varepsilon$-machine is the Fibonacci tiling]\label{ex:fib}
Let $\nu$ be the Markov measure forbidding $11$, with $\nu(0\mid0)=p$. Its classes are ``last symbol $0$'' ($A_n$) and ``last symbol $1$'' ($B_n$), of sizes $F_{n+1},F_n$, with incidence $\begin{psmallmatrix}1&1\\1&0\end{psmallmatrix}$: the Fibonacci Bratteli diagram (verified numerically to depth 14). Then $K_0(B)=\Z^2$ with the Perron--Frobenius order, and the unique trace is the Parry measure (gap labels $\Z+\Z\phi$, normalized). The KMS$_\beta$ state is unique with $x^{(n)}=r(\beta)^{-n}v_\beta$, where $r(\beta)$ is the Perron root of $\begin{psmallmatrix}p^\beta&(1-p)^\beta\\1&0\end{psmallmatrix}$: $r(0)=\phi$, $r(1)=1$, and $\log r$ is a pressure. After telescoping and a standard proper order, this stationary diagram is the Bratteli--Vershik model of the Fibonacci substitution (Durand--Host--Skau). The continuation algebra is the AF core of the Fibonacci tiling algebra, and $H^1_{PV}\cong\Z^2\cong K_0(B)$.
\end{example}

\subsection{One-dimensional tiling realization}

\begin{theorem}[realization]\label{thm:realize}
Suppose the continuation diagram (of $\nu$ or of a relaxation $\tilde\nu$) is simple, carries a proper order, has finite rank ($|P_n|\le N$ after telescoping) and topological rank $\ge2$. Then:
\begin{enumerate}
\item The Bratteli--Vershik map $\theta$ is minimal and, by Downarowicz--Maass, expansive, hence conjugate to a minimal subshift.
\item Its suspension is the hull of an aperiodic, repetitive, FLC tiling $\mathcal T_\nu$ of $\R$.
\item Its orbit relation is $G_B$ with the maximal and minimal tails merged (Herman--Putnam--Skau).
\item $H^0_{PV}=\Z$ and $H^1_{PV}\cong K^0(X_B,\theta)\cong K_0(B)$ as ordered groups (HPS, together with Putnam's AF subalgebra).
\item Invariant measures are the traces of $B$, so there are at most $N$ ergodic measures.
\item The neural state is a KMS$_1$ quasi-invariant measure with locally constant cocycle, not a BBG positive cycle unless it is a trace.
\end{enumerate}
Stationary finite $\varepsilon$-machines give substitution tilings (Durand--Host--Skau). The fair binary case ($|P_n|=1$) gives the dyadic odometer, which is not expansive; its FLC realizations are almost one-to-one Toeplitz extensions (e.g.\ period doubling), consistent with Theorem~\ref{thm:odometer}.
\end{theorem}

\begin{proposition}[diffusion--translation commutator from continuation complexity]\label{prop:complexity}
In the Bratteli--Vershik model, the level-$n$ Kakutani--Rokhlin towers are the classes $C\in P_n$, with levels the prefix cylinders $[p]$, $p\in C$, of height $|C|$. For an invariant trace $t$ with class masses $t_C$ (the mass of one prefix in $C$),
\[\|[E_n,U]f\|_2\le2\|f\|_\infty\Big(\sum_{C\in P_n^*}t_C\Big)^{1/2}\le2\|f\|_\infty\big(|P_n|\max_Ct_C\big)^{1/2},\]
with $P_n^*$ the classes whose top maps into at least two bases. For Fibonacci this is $O(\phi^{-n/2})$. Bounded continuation complexity makes the transverse diffusion and the tiling translation nearly commute at fine levels.
\end{proposition}

\begin{proof}
The Vershik successor changes only edges at levels $\le n$ unless $p$ is maximal in $C$, so it maps $[p]$ onto $[\mathrm{succ}_C(p)]$. Then apply Proposition~\ref{prop:kr}.
\end{proof}

\begin{problem}[neural rate--distortion]
Define $R(N)=\inf\{I_\nu(P;Z\mid C)\}$ over refinement-coherent relaxations with at most $N$ classes per level. By the continuation note's bound, readouts move by $\le\sqrt{R(N)/2}$. For such relaxations, the theorems above give at most $N$ extremal KMS states, Birkhoff uniqueness, a 1-D FLC realization and commutator decay. The central open problem becomes quantitative: \emph{for which source representations of the Gaussian direction is $R(N)$ small for the gains of a random ReLU network?}
\end{problem}

\section{Bridges back to the network}\label{sec:net}

\textbf{Challenge facts.} The current WhestBench Phase~2 setting, from the starter kit, is: width $1024$, depth $16$, He-normal weights $N(0,2/n)$, no biases, ReLU after every layer, inputs $N(0,I)$; the score is final-layer MSE of post-ReLU means times $\max(0.1,C/B)$; the budget is $B=2^{41}\approx2.2\times10^{12}$ FLOPs per MLP (not $241{,}241$; the digits were ``$2^{41}$''), about $6.5\times10^4$ forward passes.

\textbf{Weight-side arrows.} Fix the upstream law of $h=h_{\ell-1}(X)$ with mean $m$ and covariance $\Sigma$, a row prior $N(0,sI)$ with $s=2/n$, and a retained subspace $V$ with projection $P$. Let $f(w)=\E_X\relu(w\cdot h)$. Conditional on $Pw$, the residual $(I-P)w$ is Gaussian and independent of it. The arrows of this ``fair'' weight hierarchy are residual rotations, and partial annealing $\E[f\mid Pw]$ is the corresponding conditional expectation.
\begin{enumerate}
\item \emph{Transfer identity:} $\E[f\mid Pw]=\E_X[(P_{\tau(X)}\relu)(Pw\cdot h)]$, with heat-smoothed ReLU and diffusion time $\tau(X)=s\|(I-P)h(X)\|^2$. Discarded weight resolution equals heat flow on the activation.
\item \emph{Exact chaos formula:} $\Var(f\mid Pw)=\sum_{k\ge1}\frac1{k!}\E_{X,X'}[\psi_k\psi_k'K^k]$, with $\psi_k=(P_\tau\relu)^{(k)}(Pw\cdot h)$ and $K=s\langle(I-P)h,(I-P)h'\rangle$ (Hermite expansion of $u\mapsto\relu(a+u\cdot h')$ in the Gaussian $u$).
\item \emph{The mean arrow is essential:} if $m\in V$, the Gaussian Poincar\'e inequality gives $\Var(f\mid Pw)\le\frac s4\|(I-P)\Sigma(I-P)\|_{op}$. If $m\notin V$, then $\Var\ge s(\E\psi_1\,\|(I-P)m\|-\tfrac12\|\Sigma\|^{1/2})_+^2=\Theta(1)$.
\item \emph{Second arrow:} the variance readout $w^\top\Sigma w$ removes $99.6\%$ of the dominant second chaos, versus $24\%$ for the row norm. Annealing the Edgeworth corrections over the residual arrows produces per-layer \emph{full-trace} cumulant contractions. This gives a derivation of the ``extra full trace'' device in ARC's cumulant propagation (Wu et al., arXiv:2605.05179), and in a one-layer test ($n=256$) it cut the Gaussian-closure error from $7.1\times10^{-6}$ to $2.2\times10^{-6}$.
\end{enumerate}

\textbf{Exact downstream identity.} If layer $\ell+1$ has i.i.d.\ $N(0,s)$ weights and $G$ is the (positively homogeneous) downstream map, then for any two upstream activation laws,
\[\E_{W_{\ell+1}}(f_j-f_j')^2=\sum_{r\ge0}\kappa_r\big\|A_r(h)-A_r(h')\big\|_F^2,\qquad A_r(h)=\E_X\big[|h|\,\theta^{\otimes r}\big],\ \theta=h/|h|,\]
with $\kappa_r=\frac{s^r}{r!}\|\E\nabla^rG(U)\|^2\ge0$ the Taylor coefficients of the downstream two-point function. This follows from homogeneity, rotation invariance and multivariate Mehler. It is the weight-side analogue of Theorem~\ref{thm:seedlen}: a reduction that is pseudorandom must match the moment orders $r$ weighted by $\kappa_r$. Entries of $\E_X h'^{\otimes r}$ depend on $r$ rows jointly, so $r$-wise independent residual rows reproduce order $r$ (up to the amplitude normalization), and a shared residual direction (1-wise) fails at $r=2$. In a test ($n=256$, $L=8$, reduction at layer 3): fresh residuals gave MSE $6.5\times10^{-4}$ (identity prediction $5.8\times10^{-4}$); rows in a random 32-dimensional subspace gave $1.5\times10^{-3}$ ($1.3\times10^{-3}$); 4-dimensional, $7.3\times10^{-3}$ ($7.6\times10^{-3}$); one shared direction, $6.0\times10^{-2}$; signal variance $0.53$.

\textbf{Depth is a critical Galton--Watson process.} $\varrho(c)=\frac c2+\frac1\pi(\sqrt{1-c^2}+c\arcsin c)=\frac1\pi+\frac c2+\frac{c^2}{2\pi}+\frac{c^4}{24\pi}+\cdots$ has nonnegative coefficients summing to $1$, so it is an offspring generating function. The infinite-width downstream spectrum is $b^{(D)}_r=P(Z_D=r)$, with $\varrho^{\circ D}=\sum_rb_r^{(D)}c^r$. The mean offspring is $\varrho'(1)=\tfrac12+\tfrac1\pi\arcsin1=1$, so He initialization is \emph{exactly} criticality. Since $\varrho(1-x)=1-x+\frac{2\sqrt2}{3\pi}x^{3/2}+O(x^{5/2})$, the offspring tail is stable with index $3/2$ (infinite variance), and $P(Z_D>0)\sim9\pi^2/(2D^2)$ (Slack; this agrees with the known ReLU edge-of-chaos rate). Numerically $P(Z_{16}>0)=0.0705$, and $D^2P(Z_D>0)=43.1$ at $D=1000$, approaching $9\pi^2/2=44.41$. Contrast with PB: there the multipliers decay doubly exponentially and the gap is $\lambda_1$, giving the $\log\log$ seed. Here criticality means no gap and only polynomial forgetting with depth.

\section{Open problems}
\begin{enumerate}
\item Close the $O(N^2/\varepsilon^2)$ versus $\Omega(N)$ gap for sparse hierarchical sparsifiers (Thm.~\ref{thm:sparse}, Prop.~\ref{prop:lower}).
\item Find KMS-symmetric sparse diffusions for neural states beyond dephasing plus Metropolis, with gap bounds in terms of local rather than global distortion.
\item Develop heat--PV theory for $d\ge2$: characterize compatible hierarchies and bound $[E_l,d_{PV}]$ by BBG branch weights.
\item Compare the SB and PY spectral sequences (PY Rem.~4.1), and find a KMS-twisted analogue for nontracial (neural) states, where gap labels are replaced by $\beta$-weighted class masses.
\item Neural rate--distortion $R(N)$: find source representations making ReLU-gain continuation laws low-complexity.
\item Use the Galton--Watson law to design depth-adaptive pseudorandom weight reductions with certified error.
\end{enumerate}

\appendix
\section{Numerical checks}\label{app:num}
Scripts are in \texttt{code/}.
\begin{itemize}
\item \texttt{arrow\_checks.py}: Weyl diagonalization of $T_t$ on 5 qubits (error $0$). Sparse generator from random tail sets with $k=40$ (empirical biases $\le0.5$): $\hat\psi/\lambda\in[0.70,1.25]\subset[1-\varepsilon,1+\varepsilon]$. Trace preserved to $10^{-13}$. Choi matrix positive. $L^2$ error $3\times10^{-2}$ to $4\times10^{-3}$ against the bound $0.37$. Renormalization identity residual $2\times10^{-16}$.
\item \texttt{modular\_check.py}: reproduces the continuation note's two-bit example. The level-1 arrow is not a modular eigenoperator; all level-2 arrows are, with eigenvalues $\nu[p]/\nu[q]$.
\item \texttt{fib\_check.py}: golden-mean classes have Fibonacci sizes and incidence $\begin{psmallmatrix}1&1\\1&0\end{psmallmatrix}$; Perron roots $r(0)=\phi$, $r(1)=1$.
\item \texttt{tables.py}, \texttt{spectrum.py}: seed-length table; Galton--Watson spectra and survival asymptotics.
\item \texttt{one\_layer.py}, \texttt{chaos\_check.py}, \texttt{second\_arrow.py}, \texttt{level2.py}, \texttt{annealed\_edgeworth.py}, \texttt{depth.py}: weight-side results of Section~\ref{sec:net}. The chaos formula matches Monte Carlo to $\approx1\%$. With the mean arrow, $\E\Var=3.5\times10^{-4}$ against $0.42$ without it.
\end{itemize}

\end{document}
